\documentclass[11pt]{ctexart}
\usepackage[a4paper,margin=1in]{geometry}
\usepackage{graphicx}
\usepackage{xcolor}
\usepackage{hyperref}
\definecolor{refgray}{gray}{0.45}
\title{\textbf{市场最新观点汇总}\\[0.4em]{\large 全球资产定价锚点正从需求侧转向供给侧，地缘政治与能源冲击重塑利率、汇率与行业格局。}}
\author{KC桌面 自动生成}
\date{260524}
\begin{document}
\maketitle
\section*{一页摘要}
\begin{itemize}
  \item 中国资产定价的锚正在从增长斜率转向结构稳定性，地缘风险溢价快速压缩，但内需疲软与出口强劲的K型分化持续加深。
  \item 美联储可能无限期按兵不动，市场对降息的定价面临系统性重估，美债长端的稳定仅为技术性喘息。
  \item 全球供应链正经历供给侧再定价，地缘冲突、AI投资与能源资本开支三重力量重塑汽车、医疗科技等行业的利润分配。
  \item 人民币国际化已跨越贸易结算阶段，但资产端深度不足成为下一阶段瓶颈，境外参与者需更完善的流动性、对冲与资产生态。
  \item 机构资金正从软件转向半导体，但对冲基金与共同基金在金融和消费板块上的方向性分歧，可能孕育超额收益机会。
\end{itemize}
\section{宏观与利率：美联储无限期按兵不动，美债长端稳定缺乏基本面支撑}
\textbf{美联储降息预期正被系统性重估，美债长端的暂时企稳更多源于外部因素缓和，而非美国自身基本面改善。}

\begin{itemize}
  \item 美联储内部鹰派共识加速形成，鸽派官员也不再主张降息，政治压力消退后政策框架可能转向无限期按兵不动。
  \item 美债长端收益率回落主要受能源价格下跌及英日利好驱动，估值模型显示5y5y远期利率仅略高于公允值，尚未进入便宜区间。
  \item 欧洲利率多头性价比更高，欧洲央行政策路径更明确，而英国国债估值从宏观角度看已具吸引力。
  \item 凸性风险仍存，抵押贷款机构对冲行为不透明，但隐含波动率反应相对温和，反弹可能修复久期缺口。
  \item 全球利率分化交易正在取代简单的降息预期，欧元区资金持续买入美债以对冲自身风险，亚洲资金则持续流出。
\end{itemize}
\begin{figure}[htbp]
\centering
\includegraphics[width=0.88\linewidth]{figures/fig_025.jpg}
\caption{Exhibit 2 - Exhibit 1: Longer term UST forwards are only modestly cheap to our macro framework's estimate of fair value 5y5y UST yield vs model fair value with \$+/-1\$ standard deviation bands}
\end{figure}
\begin{figure}[htbp]
\centering
\includegraphics[width=0.88\linewidth]{figures/fig_026.jpg}
\caption{Exhibit 2 - Exhibit 2: Implied vol has lagged the magnitude of the rate shift 1m trailing range on 10y swap vs 1m10y implied vol}
\end{figure}
\begin{figure}[htbp]
\centering
\includegraphics[width=0.88\linewidth]{figures/fig_027.jpg}
\caption{Exhibit 4 - Exhibit 3: The reduction in levered fund shorts in Treasury futures may in part reflect the diminished carry in Treasury basis Leveraged fund net UST futures position (notional, adjusted for valuation changes)}
\end{figure}
\begin{figure}[htbp]
\centering
\includegraphics[width=0.88\linewidth]{figures/fig_028.jpg}
\caption{Exhibit 4 - Exhibit 4: Bunds have been a receiver rather than provider of bearish shocks G4 yields decomposition following our spillover framework}
\end{figure}
\begin{figure}[htbp]
\centering
\includegraphics[width=0.88\linewidth]{figures/fig_030.jpg}
\caption{Exhibit 6 - Exhibit 6: UK 5y5y levels look stretched from a macro perspective}
\end{figure}
{\small\color{refgray}\textbf{References:} [R006] 美债长端的“稳定”是喘息，不是反转; [R010] 市场真正低估的不是降息时点，而是美联储“无限期按兵不动”的定价}

\section{AI/算力：投资瓶颈从技术转向供给，欧洲电网与数据中心成关键约束}
\textbf{AI投资的真正瓶颈已从技术突破转向电力与数据中心等供给侧基础设施，欧洲公用事业会议和基础设施会议均聚焦于此。}

\begin{itemize}
  \item 欧洲数据中心建设成为硬需求，核心瓶颈在于电力和电网容量，利好电气化相关标的及表后解决方案提供商。
  \item AI应用的经济性仍存争议，多数企业目前使用AI尚未盈利，Token增长速度和使用盲目性是关键观察点。
  \item 机构资金正从软件板块系统性转向半导体，对冲基金半导体持仓创历史新高，软件持仓降至2019年以来最低。
  \item 共同基金剔除微软后软件低配幅度为2012年以来最大，两者均在Q2净减持微软，共识轮动方向明确。
\end{itemize}
\begin{figure}[htbp]
\centering
\includegraphics[width=0.88\linewidth]{figures/fig_045.jpg}
\caption{Exhibit 5 - Exhibit 5: Mutual funds and hedge funds have rotated from Software to Semis}
\end{figure}
\begin{figure}[htbp]
\centering
\includegraphics[width=0.88\linewidth]{figures/fig_044.jpg}
\caption{Exhibit 4 - Exhibit 4: Hedge fund vs. mutual fund sector positions Sector tilts of large-cap mutual funds and hedge funds}
\end{figure}
{\small\color{refgray}\textbf{References:} [R004] 市场真正低估的不是需求，而是供给侧的再定价; [R012] 市场共识的裂缝：机构正在从软件转向半导体，但真正的分歧藏在金融与消费里}

\section{能源与大宗：供给侧再定价深化，地缘冲突成本向2027年累积}
\textbf{全球能源与大宗商品市场正经历供给侧结构性再定价，地缘冲突对汽车供应链和欧洲通胀的影响比市场预期的更持久。}

\begin{itemize}
  \item 全球汽车供应链面临地缘冲突隐性成本上升、非汽车业务争夺产能、中国产能外溢三重供给冲击，2027年利润率压力将显著大于当前预期。
  \item 欧洲天然气冲击较2022年温和，但石油和成品油价格预期路径显示更强粘性，远期曲线比天然气更陡。
  \item 石油资本开支勘探周期已启动，地震勘探服务商率先受益，行业整合后活跃船舶数量大幅下降。
  \item 能源流动受限正重塑全球贸易条件，美元升值压力持续，市场可能低估了能源溢价的持久性。
  \item 中国汽车出口因全球油气价格上涨而加速，4月出口同比暴增84.1\%，但出口量仅占总量19\%，难以弥补国内需求缺口。
\end{itemize}
\begin{figure}[htbp]
\centering
\includegraphics[width=0.88\linewidth]{figures/fig_017.jpg}
\caption{Exhibit 1 - Exhibit 1: Current Shock is Seen as More Contained for European Natural Gas Than in 2022; Oil Shock Expected to Be More Persistent}
\end{figure}
\begin{figure}[htbp]
\centering
\includegraphics[width=0.88\linewidth]{figures/fig_018.jpg}
\caption{Exhibit 2 - Exhibit 2: Spot Gasoline Prices and Electricity Forward Prices Have Risen Since the Onset of the Conflict But Are Below Their Recent Peaks}
\end{figure}
\begin{figure}[htbp]
\centering
\includegraphics[width=0.88\linewidth]{figures/fig_023.jpg}
\caption{Exhibit 7 - Exhibit 7: We Have Raised Our End-2026 Inflation Forecast By Around 2.0pp in the Euro Area and 1.7pp in the UK}
\end{figure}
\begin{figure}[htbp]
\centering
\includegraphics[width=0.88\linewidth]{figures/fig_031.jpg}
\caption{Exhibit 1 - Exhibit 1: The net effect of the AI boom and higher-for-longer energy prices leaves the US looking like a relative outperformer}
\end{figure}
\begin{figure}[htbp]
\centering
\includegraphics[width=0.88\linewidth]{figures/fig_035.jpg}
\caption{Exhibit 5 - Exhibit 5: Long NOK/SEK has been on par with long July Brent for vol-adjusted returns in the energy shock Vol-Adjusted Move Since Feb 27}
\end{figure}
{\small\color{refgray}\textbf{References:} [R002] 市场真正低估的不是需求，而是供给侧的再定价; [R005] 市场真正低估的不是地缘风险，而是欧洲通胀的二次定价; [R007] 美元正在被重新定价，但市场可能低估了“能源溢价”的持久性; [R009] 中国车市真正的分水岭：不是需求疲软，而是内外市场正在走向“脱钩”}

\section{地产与消费：中国内需疲软是政策红利退潮后的真实水位，消费信心修复需名义工资增长}
\textbf{中国国内消费与地产的结构性疲软并非周期性波动，而是政策刺激退潮后的真实需求水位，消费反弹依赖名义工资增长。}

\begin{itemize}
  \item 中国汽车国内零售销量5月预计同比下滑21.6\%，主要受以旧换新补贴收缩和新能源车购置税恢复至5\%的叠加影响。
  \item 居民存款积累约50-60万亿但不敢消费，核心原因包括失业担忧、工资增长弱和地产财富缩水。
  \item AI对服务业岗位的替代效应正在加剧消费疲软，结构性障碍短期内难以消除。
  \item 新经济（AI设备、新能源）增速强劲，AI相关设备出口同比接近翻倍，但旧经济（消费、地产、服务业）持续拖累。
  \item 国内与出口市场的增长裂痕在2026年将持续扩大，汽车行业占GDP的3.2\%，内需下滑对整体经济拖累不容小觑。
\end{itemize}
{\small\color{refgray}\textbf{References:} [R001] 中国真正的叙事切换：不是增长，是资产定价的锚在变; [R009] 中国车市真正的分水岭：不是需求疲软，而是内外市场正在走向“脱钩”}

\section{区域市场：欧洲资金回流美国但结构分化，新兴市场内部换仓明显}
\textbf{全球资本流动呈现高度分化，欧洲资金重新买入美债，亚洲资金持续流出，新兴市场内部从A股转向韩国。}

\begin{itemize}
  \item 截至5月20日当周，全球股票基金净流入约24亿美元，债券基金净流入约32亿美元，主要靠国债和综合型债券基金支撑。
  \item 欧元区年初以来累计流入美国国债基金超110亿美元，亚洲地区净流出超50亿美元，剪刀差持续扩大。
  \item 科技板块成为资金最爱，金融板块遭遇最大净流出，工业、能源、基础设施板块年初至今累计流入最多。
  \item 全球新兴市场基准基金和A股基金出现净流出，但韩国基金重新获得净流入，新兴市场本币债券基金净流入。
  \item 欧洲医疗科技板块估值可能已触及周期底部，盈利下修速度收窄，但市场需要看到多个季度企稳信号才会重新入场。
\end{itemize}
\begin{figure}[htbp]
\centering
\includegraphics[width=0.88\linewidth]{figures/fig_039.jpg}
\caption{Figure 1 - Figure 1: DB EU MedTech \& Life Sciences coverage BBG earnings revisions (indexed)}
\end{figure}
\begin{figure}[htbp]
\centering
\includegraphics[width=0.88\linewidth]{figures/fig_040.jpg}
\caption{Figure 2 - Figure 2: DB EU MedTech \& Life Sciences next FY P/E estimates}
\end{figure}
\begin{figure}[htbp]
\centering
\includegraphics[width=0.88\linewidth]{figures/fig_051.jpg}
\caption{Exhibit 14 - Exhibit 14: Recent mutual fund and ETF flows}
\end{figure}
{\small\color{refgray}\textbf{References:} [R008] 欧洲正在重新买入美国资产，但真正的结构性分歧藏在资金流向的细节里; [R011] 欧洲医疗科技板块的估值底部已经出现，但市场需要的不只是一个好季度}

\section{风险偏好：机构从软件转向半导体，对冲基金与共同基金在金融和消费上方向相反}
\textbf{市场共识正在从软件转向半导体，但真正的分歧在于对冲基金与共同基金在金融和可选消费板块上的方向性差异。}

\begin{itemize}
  \item 对冲基金年初至今回报约7\%，收益主要来自Beta和长期多头Alpha，重仓篮子回报13\%，做空篮子回报27\%。
  \item 共同基金仅26\%跑赢基准，远低于32\%的历史均值，传统选股框架有效性下降。
  \item 对冲基金净杠杆率已升至近五年第85百分位，总杠杆率仍处历史高位；共同基金现金占比小幅提升至1.4\%。
  \item 标普500中位数个股空头比例占市值3\%，为2011年以来最高，做空情绪升温。
  \item 对冲基金超配工业、低配信息技术，共同基金超配金融、低配可选消费，两者方向截然相反。
\end{itemize}
\begin{figure}[htbp]
\centering
\includegraphics[width=0.88\linewidth]{figures/fig_041.jpg}
\caption{Exhibit 1 - Exhibit 1: Equity hedge funds have returned \$7\textbackslash{}\%\$ YTD}
\end{figure}
\begin{figure}[htbp]
\centering
\includegraphics[width=0.88\linewidth]{figures/fig_042.jpg}
\caption{Exhibit 2 - Exhibit 2: \$30\textbackslash{}\%\$ of large-cap mutual funds have outperformed benchmarks Share of mutual funds outperforming benchmarks}
\end{figure}
\begin{figure}[htbp]
\centering
\includegraphics[width=0.88\linewidth]{figures/fig_043.jpg}
\caption{Exhibit 3 - Exhibit 3: Hedge fund and mutual fund exposure to equities hedge fund data as of May 21, 2026; mutual fund data as of March 31, 2026}
\end{figure}
\begin{figure}[htbp]
\centering
\includegraphics[width=0.88\linewidth]{figures/fig_044.jpg}
\caption{Exhibit 4 - Exhibit 4: Hedge fund vs. mutual fund sector positions Sector tilts of large-cap mutual funds and hedge funds}
\end{figure}
\begin{figure}[htbp]
\centering
\includegraphics[width=0.88\linewidth]{figures/fig_045.jpg}
\caption{Exhibit 5 - Exhibit 5: Mutual funds and hedge funds have rotated from Software to Semis}
\end{figure}
{\small\color{refgray}\textbf{References:} [R012] 市场共识的裂缝：机构正在从软件转向半导体，但真正的分歧藏在金融与消费里}

\section{人民币国际化：贸易结算领跑，但资产端深度不足成为下一阶段瓶颈}
\textbf{人民币跨境使用已从贸易结算转向投资驱动，但境外参与者获取流动性、对冲风险和寻找优质资产的能力仍远落后于交易量增长。}

\begin{itemize}
  \item 人民币跨境交易从2017年的9万亿元跃升至2024年的64万亿元，资本和金融账户交易占比已达75\%。
  \item 债券投资占跨境人民币支付的46\%，货物贸易仅占19\%，投资驱动成为主要动力。
  \item 人民币在全球支付份额约3-4\%，官方储备份额约2\%，国际债券份额仅0.9\%，远低于中国占全球GDP约19\%的体量。
  \item 离岸人民币债券发行在2026年前四个月几乎翻倍，但资产端配套仍不完善，CNH-CNY融资利差虽已收窄但仍存在。
  \item 境外投资者增持人民币股票但减持债券，外汇衍生品市场人民币参与度远高于利率衍生品，利率风险管理工具不足。
\end{itemize}
\begin{figure}[htbp]
\centering
\includegraphics[width=0.88\linewidth]{figures/fig_001.jpg}
\caption{Exhibit 1 - Exhibit 1: Average daily RMB transaction volume of CIPS spiked in March Average daily RMB transaction volume of CIPS}
\end{figure}
\begin{figure}[htbp]
\centering
\includegraphics[width=0.88\linewidth]{figures/fig_002.jpg}
\caption{Exhibit 2 - Exhibit 2: Dim sum bond issuance nearly doubled from a year ago in Jan-Apr 2026 Dim Sum Gross Issuance (cumulative)}
\end{figure}
\begin{figure}[htbp]
\centering
\includegraphics[width=0.88\linewidth]{figures/fig_003.jpg}
\caption{Exhibit 3 - Exhibit 3: RMB international use has improved, but still trails China's global GDP and trade shares}
\end{figure}
\begin{figure}[htbp]
\centering
\includegraphics[width=0.88\linewidth]{figures/fig_005.jpg}
\caption{Exhibit 5 - Exhibit 5: Financial flows have become the main driver of cross-border RMB transactions Cross-border RMB transaction volume}
\end{figure}
\begin{figure}[htbp]
\centering
\includegraphics[width=0.88\linewidth]{figures/fig_006.jpg}
\caption{Exhibit 6 - Exhibit 6: Bond investment accounted for nearly half of cross-border RMB transactions in 2024 Cross-border RMB transactions by usage (2024)}
\end{figure}
\begin{figure}[htbp]
\centering
\includegraphics[width=0.88\linewidth]{figures/fig_008.jpg}
\caption{Exhibit 8 - Exhibit 8: CNH-CNY funding spreads have narrowed and stabilized since April 2025}
\end{figure}
\begin{figure}[htbp]
\centering
\includegraphics[width=0.88\linewidth]{figures/fig_011.jpg}
\caption{Exhibit 11 - Exhibit 11: RMB's footprint in global derivatives markets is much larger in FX than in rates Share of RMB in global FX/rates derivatives turnover}
\end{figure}
\begin{figure}[htbp]
\centering
\includegraphics[width=0.88\linewidth]{figures/fig_016.jpg}
\caption{Exhibit 16 - Exhibit 16: Foreign investors have increased their holdings of RMB equities but reduced holdings of RMB bonds over the past few months}
\end{figure}
{\small\color{refgray}\textbf{References:} [R003] 人民币国际化真正的瓶颈不在贸易结算，而在资产端的深度}

\section*{结语}
综合来看，市场正处于定价锚切换的关键阶段：供给侧因素（地缘冲突、能源资本开支、AI基础设施）正在取代传统的需求侧逻辑，成为驱动资产价格的核心变量。投资者需警惕旧框架失效带来的风险，同时关注共识裂缝中孕育的结构性机会。
\vfill
{\small\color{refgray}Personal reading notes and learning share only. Not investment advice.}
\end{document}
